\documentclass[11pt,a4paper]{article}

% ------------------------------------------------------------------------------
% Packages
% ------------------------------------------------------------------------------
\usepackage[utf8]{inputenc}
\usepackage[T1]{fontenc}
\usepackage[french]{babel}
\usepackage{lmodern}
\usepackage[a4paper,margin=2.2cm]{geometry}
\usepackage{xcolor}
\usepackage{booktabs}
\usepackage{multirow}
\usepackage{array}
\usepackage{float}
\usepackage{caption}
\usepackage{listings}
\usepackage{tikz}
\usetikzlibrary{babel,arrows.meta}
\usepackage{hyperref}

\hypersetup{colorlinks=true, linkcolor=bleu, urlcolor=bleu}

% ------------------------------------------------------------------------------
% Couleurs
% ------------------------------------------------------------------------------
\definecolor{bleu}{HTML}{1E3A8A}
\definecolor{chemin}{HTML}{2563EB}
\definecolor{cheminfond}{HTML}{DBEAFE}
\definecolor{vert}{HTML}{16A34A}
\definecolor{vertfond}{HTML}{DCFCE7}
\definecolor{orange}{HTML}{EA580C}
\definecolor{orangefond}{HTML}{FFEDD5}
\definecolor{rouge}{HTML}{DC2626}
\definecolor{gris}{HTML}{64748B}
\definecolor{codefond}{HTML}{F8FAFC}
\definecolor{codecadre}{HTML}{CBD5E1}

% ------------------------------------------------------------------------------
% Style du code C++
% ------------------------------------------------------------------------------
\lstset{
  language=C++,
  basicstyle=\ttfamily\footnotesize,
  keywordstyle=\color{chemin}\bfseries,
  commentstyle=\color{gris}\itshape,
  stringstyle=\color{vert},
  backgroundcolor=\color{codefond},
  frame=single,
  rulecolor=\color{codecadre},
  breaklines=true,
  breakatwhitespace=true,
  columns=fullflexible,
  keepspaces=true,
  showstringspaces=false,
  tabsize=4,
  numbers=left,
  numberstyle=\tiny\color{gris},
  captionpos=b,
}
\renewcommand{\lstlistingname}{Code}

% ------------------------------------------------------------------------------
% Styles TikZ
% ------------------------------------------------------------------------------
\tikzset{
  noeud/.style={circle, draw=black!75, thick, fill=white, minimum size=8mm,
                inner sep=0pt, font=\small\bfseries},
  actif/.style={noeud, draw=orange, fill=orangefond},
  neuf/.style={noeud, draw=vert, fill=vertfond},
  arete/.style={draw=black!45, thick},
  echange/.style={<->, >={Stealth[length=2.5mm]}, rouge, very thick},
  case/.style={rectangle, draw=black!60, thick, minimum width=8mm, minimum height=7mm,
               inner sep=0pt, font=\small\bfseries},
  noeudcase/.style={fill=white},
  actifcase/.style={draw=orange, fill=orangefond},
  neufcase/.style={draw=vert, fill=vertfond},
  indice/.style={font=\scriptsize, text=gris},
}

\setlength{\parindent}{0pt}
\setlength{\parskip}{5pt}
\renewcommand{\arraystretch}{1.15}

\title{\textbf{Tas min}\\[4pt]
       \large Insertions et extraction du min de 200\,000 à 10\,000\,000 valeurs, comparaison avec \texttt{std::priority\_queue}}
\author{Zakariae Bakkari\\
        \small Programmation Objet Avancée et Design Patterns\\
        \small\url{https://github.com/zakariae-bakkari/cpp_ilisi2}}
\date{30 octobre 2026}

\begin{document}

\maketitle

% ==============================================================================
\section{Objectif}
% ==============================================================================

Implémenter un \textbf{tas min} en C++ (\texttt{src/12\_tas\_min.cpp}), y insérer les
valeurs d'un fichier (200\,000, 1\,000\,000 puis 10\,000\,000 valeurs), puis mesurer le temps
et le nombre d'échanges de trois opérations sur le tas construit :
\textbf{extraire le minimum une fois} (retirer la racine, puis réorganiser le tas),
\textbf{insérer une valeur qui reste à la fin} (aucune réorganisation) et
\textbf{insérer une valeur qui remonte jusqu'à la racine} (vérifiée à chaque niveau).
Puis comparer avec le tas min de la STL, \texttt{std::priority\_queue} avec
\texttt{greater<int>}.

\textbf{Code source :}
\begin{itemize}
  \item notre tas (option \textbf{13} du menu \texttt{./run.sh}, ou \texttt{make run\_12}) :\\
        \url{https://github.com/zakariae-bakkari/cpp_ilisi2/blob/main/iterators/src/12_tas_min.cpp}
  \item comparaison notre tas / \texttt{std::priority\_queue} (option \textbf{14}, ou \texttt{make run\_13}) :\\
        \url{https://github.com/zakariae-bakkari/cpp_ilisi2/blob/main/iterators/src/13_tas_min_priority_queue.cpp}
\end{itemize}

% ==============================================================================
\section{Un arbre rangé dans un tableau}
% ==============================================================================

Un tas min est un arbre binaire \textbf{complet} (rempli niveau par niveau, de gauche à
droite) qui respecte une seule règle : \textbf{chaque parent est plus petit que ses
enfants}. Le minimum est donc toujours à la racine. C'est le tas max avec la comparaison
inversée.

Comme l'arbre est complet, on le range dans un simple \texttt{vector<int>}, sans pointeurs.
Pour la case $i$ :

\begin{center}
\begin{tabular}{l l}
parent        & $(i - 1) / 2$ \\
enfant gauche & $2i + 1$ \\
enfant droit  & $2i + 2$ \\
\end{tabular}
\end{center}

\begin{figure}[H]
\centering
\begin{minipage}{0.5\textwidth}\centering
\begin{minipage}[t]{0.32\textwidth}\centering
\begin{tikzpicture}[scale=0.62, transform shape]
  \draw[arete] (0,0) -- (-1.6,-1.1);
  \draw[arete] (0,0) -- (1.6,-1.1);
  \draw[arete] (-1.6,-1.1) -- (-2.4,-2.2);
  \draw[arete] (-1.6,-1.1) -- (-0.8,-2.2);
  \draw[arete] (1.6,-1.1) -- (0.8,-2.2);
  \node[noeud] (n0) at (0,0) {1};
  \node[noeud] (n1) at (-1.6,-1.1) {4};
  \node[noeud] (n2) at (1.6,-1.1) {3};
  \node[noeud] (n3) at (-2.4,-2.2) {15};
  \node[noeud] (n4) at (-0.8,-2.2) {20};
  \node[noeud] (n5) at (0.8,-2.2) {10};
  \node[case, noeudcase] at (-2.00,-3.4) {1};
  \node[indice] at (-2.00,-3.95) {0};
  \node[case, noeudcase] at (-1.20,-3.4) {4};
  \node[indice] at (-1.20,-3.95) {1};
  \node[case, noeudcase] at (-0.40,-3.4) {3};
  \node[indice] at (-0.40,-3.95) {2};
  \node[case, noeudcase] at (0.40,-3.4) {15};
  \node[indice] at (0.40,-3.95) {3};
  \node[case, noeudcase] at (1.20,-3.4) {20};
  \node[indice] at (1.20,-3.95) {4};
  \node[case, noeudcase] at (2.00,-3.4) {10};
  \node[indice] at (2.00,-3.95) {5};
\end{tikzpicture}\\[2pt]
\small 
\end{minipage}
\end{minipage}
\caption{Le même tas vu comme un arbre et comme le tableau réellement stocké (indices en gris). Le parent de la case 5 (10) est la case $(5-1)/2 = 2$ (3).}
\end{figure}

% ==============================================================================
\section{Insertion : on ajoute en bas, puis on remonte}
% ==============================================================================

\begin{lstlisting}[caption={Insertion (\texttt{inserer} et \texttt{monter}).}]
// Insertion : ajoute la valeur dans la premiere case libre (a la fin),
// puis la fait remonter. O(log n).
void inserer(int val) {
    tas.push_back(val);
    monter(tas.size() - 1);
}

// Fait remonter la case i : tant qu'elle est plus petite que son parent,
// on les echange. S'arrete a la racine ou quand le parent est plus petit.
void monter(int i) {
    while (i > 0 && tas[(i - 1) / 2] > tas[i]) {
        swap(tas[i], tas[(i - 1) / 2]);
        echanges++;
        i = (i - 1) / 2;
    }
}
\end{lstlisting}

La nouvelle valeur est placée dans la première case libre (à la fin du tableau). Tant
qu'elle est plus \textbf{petite} que son parent, on les \textbf{échange} : elle remonte d'un
niveau. Elle s'arrête quand son parent est plus petit, ou à la racine.

Exemple : on insère dans l'ordre $4, 15, 10, 3, 20, 1$. Voici l'insertion de \textbf{3}
dans le tas $[4, 15, 10]$ (en vert la valeur qui remonte, en orange la case comparée, en
rouge l'échange) :

\begin{figure}[H]
\centering
\begin{minipage}[t]{0.32\textwidth}\centering
\begin{tikzpicture}[scale=0.62, transform shape]
  \draw[arete] (0,0) -- (-1.6,-1.1);
  \draw[arete] (0,0) -- (1.6,-1.1);
  \draw[arete] (-1.6,-1.1) -- (-2.4,-2.2);
  \node[noeud] (n0) at (0,0) {4};
  \node[noeud] (n1) at (-1.6,-1.1) {15};
  \node[noeud] (n2) at (1.6,-1.1) {10};
  \node[neuf] (n3) at (-2.4,-2.2) {3};
  \node[case, noeudcase] at (-1.20,-3.4) {4};
  \node[indice] at (-1.20,-3.95) {0};
  \node[case, noeudcase] at (-0.40,-3.4) {15};
  \node[indice] at (-0.40,-3.95) {1};
  \node[case, noeudcase] at (0.40,-3.4) {10};
  \node[indice] at (0.40,-3.95) {2};
  \node[case, neufcase] at (1.20,-3.4) {3};
  \node[indice] at (1.20,-3.95) {3};
\end{tikzpicture}\\[2pt]
\small \textbf{1.} 3 est ajouté à la fin (case 3)
\end{minipage}
\hfill
\begin{minipage}[t]{0.32\textwidth}\centering
\begin{tikzpicture}[scale=0.62, transform shape]
  \draw[arete] (0,0) -- (-1.6,-1.1);
  \draw[arete] (0,0) -- (1.6,-1.1);
  \draw[arete] (-1.6,-1.1) -- (-2.4,-2.2);
  \node[noeud] (n0) at (0,0) {4};
  \node[actif] (n1) at (-1.6,-1.1) {15};
  \node[noeud] (n2) at (1.6,-1.1) {10};
  \node[neuf] (n3) at (-2.4,-2.2) {3};
  \draw[echange] (n3) to[bend right=35] (n1);
  \node[case, noeudcase] at (-1.20,-3.4) {4};
  \node[indice] at (-1.20,-3.95) {0};
  \node[case, actifcase] at (-0.40,-3.4) {15};
  \node[indice] at (-0.40,-3.95) {1};
  \node[case, noeudcase] at (0.40,-3.4) {10};
  \node[indice] at (0.40,-3.95) {2};
  \node[case, neufcase] at (1.20,-3.4) {3};
  \node[indice] at (1.20,-3.95) {3};
\end{tikzpicture}\\[2pt]
\small \textbf{2.} parent = case $(3-1)/2 = 1$ : $15 > 3$, on échange
\end{minipage}
\hfill
\begin{minipage}[t]{0.32\textwidth}\centering
\begin{tikzpicture}[scale=0.62, transform shape]
  \draw[arete] (0,0) -- (-1.6,-1.1);
  \draw[arete] (0,0) -- (1.6,-1.1);
  \draw[arete] (-1.6,-1.1) -- (-2.4,-2.2);
  \node[actif] (n0) at (0,0) {4};
  \node[neuf] (n1) at (-1.6,-1.1) {3};
  \node[noeud] (n2) at (1.6,-1.1) {10};
  \node[noeud] (n3) at (-2.4,-2.2) {15};
  \draw[echange] (n1) to[bend right=35] (n0);
  \node[case, actifcase] at (-1.20,-3.4) {4};
  \node[indice] at (-1.20,-3.95) {0};
  \node[case, neufcase] at (-0.40,-3.4) {3};
  \node[indice] at (-0.40,-3.95) {1};
  \node[case, noeudcase] at (0.40,-3.4) {10};
  \node[indice] at (0.40,-3.95) {2};
  \node[case, noeudcase] at (1.20,-3.4) {15};
  \node[indice] at (1.20,-3.95) {3};
\end{tikzpicture}\\[2pt]
\small \textbf{3.} parent = case 0 : $4 > 3$, on échange
\end{minipage}

\vspace{6mm}

\begin{minipage}[t]{0.32\textwidth}\centering
\begin{tikzpicture}[scale=0.62, transform shape]
  \draw[arete] (0,0) -- (-1.6,-1.1);
  \draw[arete] (0,0) -- (1.6,-1.1);
  \draw[arete] (-1.6,-1.1) -- (-2.4,-2.2);
  \node[neuf] (n0) at (0,0) {3};
  \node[noeud] (n1) at (-1.6,-1.1) {4};
  \node[noeud] (n2) at (1.6,-1.1) {10};
  \node[noeud] (n3) at (-2.4,-2.2) {15};
  \node[case, neufcase] at (-1.20,-3.4) {3};
  \node[indice] at (-1.20,-3.95) {0};
  \node[case, noeudcase] at (-0.40,-3.4) {4};
  \node[indice] at (-0.40,-3.95) {1};
  \node[case, noeudcase] at (0.40,-3.4) {10};
  \node[indice] at (0.40,-3.95) {2};
  \node[case, noeudcase] at (1.20,-3.4) {15};
  \node[indice] at (1.20,-3.95) {3};
\end{tikzpicture}\\[2pt]
\small \textbf{4.} 3 est à la racine : on s'arrête
\end{minipage}
\hfill
\begin{minipage}[t]{0.32\textwidth}\centering
\begin{tikzpicture}[scale=0.62, transform shape]
  \draw[arete] (0,0) -- (-1.6,-1.1);
  \draw[arete] (0,0) -- (1.6,-1.1);
  \draw[arete] (-1.6,-1.1) -- (-2.4,-2.2);
  \draw[arete] (-1.6,-1.1) -- (-0.8,-2.2);
  \node[noeud] (n0) at (0,0) {3};
  \node[actif] (n1) at (-1.6,-1.1) {4};
  \node[noeud] (n2) at (1.6,-1.1) {10};
  \node[noeud] (n3) at (-2.4,-2.2) {15};
  \node[neuf] (n4) at (-0.8,-2.2) {20};
  \node[case, noeudcase] at (-1.60,-3.4) {3};
  \node[indice] at (-1.60,-3.95) {0};
  \node[case, actifcase] at (-0.80,-3.4) {4};
  \node[indice] at (-0.80,-3.95) {1};
  \node[case, noeudcase] at (0.00,-3.4) {10};
  \node[indice] at (0.00,-3.95) {2};
  \node[case, noeudcase] at (0.80,-3.4) {15};
  \node[indice] at (0.80,-3.95) {3};
  \node[case, neufcase] at (1.60,-3.4) {20};
  \node[indice] at (1.60,-3.95) {4};
\end{tikzpicture}\\[2pt]
\small \textbf{Insérer 20} : parent 4 $< 20$, aucun échange
\end{minipage}
\hfill
\begin{minipage}[t]{0.32\textwidth}\centering
\begin{tikzpicture}[scale=0.62, transform shape]
  \draw[arete] (0,0) -- (-1.6,-1.1);
  \draw[arete] (0,0) -- (1.6,-1.1);
  \draw[arete] (-1.6,-1.1) -- (-2.4,-2.2);
  \draw[arete] (-1.6,-1.1) -- (-0.8,-2.2);
  \draw[arete] (1.6,-1.1) -- (0.8,-2.2);
  \node[neuf] (n0) at (0,0) {1};
  \node[noeud] (n1) at (-1.6,-1.1) {4};
  \node[noeud] (n2) at (1.6,-1.1) {3};
  \node[noeud] (n3) at (-2.4,-2.2) {15};
  \node[noeud] (n4) at (-0.8,-2.2) {20};
  \node[noeud] (n5) at (0.8,-2.2) {10};
  \node[case, neufcase] at (-2.00,-3.4) {1};
  \node[indice] at (-2.00,-3.95) {0};
  \node[case, noeudcase] at (-1.20,-3.4) {4};
  \node[indice] at (-1.20,-3.95) {1};
  \node[case, noeudcase] at (-0.40,-3.4) {3};
  \node[indice] at (-0.40,-3.95) {2};
  \node[case, noeudcase] at (0.40,-3.4) {15};
  \node[indice] at (0.40,-3.95) {3};
  \node[case, noeudcase] at (1.20,-3.4) {20};
  \node[indice] at (1.20,-3.95) {4};
  \node[case, noeudcase] at (2.00,-3.4) {10};
  \node[indice] at (2.00,-3.95) {5};
\end{tikzpicture}\\[2pt]
\small \textbf{Insérer 1} : 2 échanges (avec 10 puis 3)
\end{minipage}
\caption{Insertion pas à pas. Le tableau final est $[1, 4, 3, 15, 20, 10]$.}
\end{figure}

Une insertion a donc un \textbf{meilleur cas} et un \textbf{pire cas}, que le programme
mesure tous les deux une fois le tas construit :
\begin{itemize}
  \item \textbf{la valeur reste à la fin} (comme 20 ci-dessus) : elle est plus grande que
        son parent, une seule comparaison suffit, \textbf{0 échange}, $O(1)$ ;
  \item \textbf{la valeur remonte jusqu'à la racine} (comme 1) : elle est plus petite que
        toutes les autres, il faut la vérifier et l'échanger à chaque niveau,
        \textbf{$\log_2 n$ échanges} (17 pour 200\,000 valeurs), $O(\log n)$.
\end{itemize}
Conséquence : insérer des valeurs \textbf{déjà triées par ordre croissant} est le meilleur
cas du tas min (chaque valeur est plus grande que toutes les précédentes et reste à la fin),
alors que c'était le pire cas du tas max.

% ==============================================================================
\section{Extraction : on retire la racine, puis on redescend}
% ==============================================================================

\begin{lstlisting}[caption={Extraction du minimum (\texttt{extraireMin} et \texttt{descendre}).}]
// Extraction : retire et renvoie la racine (le min). Le dernier element
// prend sa place, puis on l'ajuste en le faisant descendre. O(log n).
int extraireMin() {
    int racine = tas[0];
    tas[0] = tas.back();
    tas.pop_back();
    if (!tas.empty()) descendre(0);
    return racine;
}

// Fait descendre la case i : on l'echange avec le plus petit de ses deux
// enfants tant qu'il est plus petit qu'elle. S'arrete quand elle est a sa
// place ou qu'elle n'a plus d'enfant.
void descendre(int i) {
    int n = tas.size();
    while (true) {
        int min = i;
        int g = 2 * i + 1;
        int d = 2 * i + 2;
        if (g < n && tas[g] < tas[min]) min = g;
        if (d < n && tas[d] < tas[min]) min = d;
        if (min == i) return;
        swap(tas[i], tas[min]);
        echanges++;
        i = min;
    }
}
\end{lstlisting}

Le minimum est la racine. Pour ne pas laisser de trou, on met le \textbf{dernier} élément
à la racine et on raccourcit le tableau. Cet élément est grand, donc la règle n'est plus
respectée : on l'\textbf{ajuste} en l'échangeant avec le plus petit de ses deux enfants,
jusqu'à ce qu'il soit plus petit que ses enfants ou qu'il n'en ait plus.

\begin{figure}[H]
\centering
\begin{minipage}[t]{0.32\textwidth}\centering
\begin{tikzpicture}[scale=0.62, transform shape]
  \draw[arete] (0,0) -- (-1.6,-1.1);
  \draw[arete] (0,0) -- (1.6,-1.1);
  \draw[arete] (-1.6,-1.1) -- (-2.4,-2.2);
  \draw[arete] (-1.6,-1.1) -- (-0.8,-2.2);
  \draw[arete] (1.6,-1.1) -- (0.8,-2.2);
  \node[actif] (n0) at (0,0) {1};
  \node[noeud] (n1) at (-1.6,-1.1) {4};
  \node[noeud] (n2) at (1.6,-1.1) {3};
  \node[noeud] (n3) at (-2.4,-2.2) {15};
  \node[noeud] (n4) at (-0.8,-2.2) {20};
  \node[neuf] (n5) at (0.8,-2.2) {10};
  \node[case, actifcase] at (-2.00,-3.4) {1};
  \node[indice] at (-2.00,-3.95) {0};
  \node[case, noeudcase] at (-1.20,-3.4) {4};
  \node[indice] at (-1.20,-3.95) {1};
  \node[case, noeudcase] at (-0.40,-3.4) {3};
  \node[indice] at (-0.40,-3.95) {2};
  \node[case, noeudcase] at (0.40,-3.4) {15};
  \node[indice] at (0.40,-3.95) {3};
  \node[case, noeudcase] at (1.20,-3.4) {20};
  \node[indice] at (1.20,-3.95) {4};
  \node[case, neufcase] at (2.00,-3.4) {10};
  \node[indice] at (2.00,-3.95) {5};
\end{tikzpicture}\\[2pt]
\small \textbf{1.} on retire la racine 1, le dernier (10) va prendre sa place
\end{minipage}
\hfill
\begin{minipage}[t]{0.32\textwidth}\centering
\begin{tikzpicture}[scale=0.62, transform shape]
  \draw[arete] (0,0) -- (-1.6,-1.1);
  \draw[arete] (0,0) -- (1.6,-1.1);
  \draw[arete] (-1.6,-1.1) -- (-2.4,-2.2);
  \draw[arete] (-1.6,-1.1) -- (-0.8,-2.2);
  \node[neuf] (n0) at (0,0) {10};
  \node[noeud] (n1) at (-1.6,-1.1) {4};
  \node[actif] (n2) at (1.6,-1.1) {3};
  \node[noeud] (n3) at (-2.4,-2.2) {15};
  \node[noeud] (n4) at (-0.8,-2.2) {20};
  \draw[echange] (n0) to[bend right=35] (n2);
  \node[case, neufcase] at (-1.60,-3.4) {10};
  \node[indice] at (-1.60,-3.95) {0};
  \node[case, noeudcase] at (-0.80,-3.4) {4};
  \node[indice] at (-0.80,-3.95) {1};
  \node[case, actifcase] at (0.00,-3.4) {3};
  \node[indice] at (0.00,-3.95) {2};
  \node[case, noeudcase] at (0.80,-3.4) {15};
  \node[indice] at (0.80,-3.95) {3};
  \node[case, noeudcase] at (1.60,-3.4) {20};
  \node[indice] at (1.60,-3.95) {4};
\end{tikzpicture}\\[2pt]
\small \textbf{2.} enfants 4 et 3 : le plus petit est 3 $< 10$, on échange
\end{minipage}
\hfill
\begin{minipage}[t]{0.32\textwidth}\centering
\begin{tikzpicture}[scale=0.62, transform shape]
  \draw[arete] (0,0) -- (-1.6,-1.1);
  \draw[arete] (0,0) -- (1.6,-1.1);
  \draw[arete] (-1.6,-1.1) -- (-2.4,-2.2);
  \draw[arete] (-1.6,-1.1) -- (-0.8,-2.2);
  \node[noeud] (n0) at (0,0) {3};
  \node[noeud] (n1) at (-1.6,-1.1) {4};
  \node[neuf] (n2) at (1.6,-1.1) {10};
  \node[noeud] (n3) at (-2.4,-2.2) {15};
  \node[noeud] (n4) at (-0.8,-2.2) {20};
  \node[case, noeudcase] at (-1.60,-3.4) {3};
  \node[indice] at (-1.60,-3.95) {0};
  \node[case, noeudcase] at (-0.80,-3.4) {4};
  \node[indice] at (-0.80,-3.95) {1};
  \node[case, neufcase] at (0.00,-3.4) {10};
  \node[indice] at (0.00,-3.95) {2};
  \node[case, noeudcase] at (0.80,-3.4) {15};
  \node[indice] at (0.80,-3.95) {3};
  \node[case, noeudcase] at (1.60,-3.4) {20};
  \node[indice] at (1.60,-3.95) {4};
\end{tikzpicture}\\[2pt]
\small \textbf{3.} 10 n'a plus d'enfant : on s'arrête, 3 est le nouveau min
\end{minipage}
\caption{Extraction du minimum du tas $[1, 4, 3, 15, 20, 10]$.}
\end{figure}

Une extraction descend au plus de $\log_2 n$ niveaux : elle coûte $O(\log n)$. Le programme
extrait le minimum \textbf{une seule fois}, après avoir inséré toutes les valeurs, puis fait
les deux insertions (reste à la fin, jusqu'à la racine).

% ==============================================================================
\section{Résultats}
% ==============================================================================

\texttt{TasMin} : notre tas écrit à la main ; \texttt{PQ} :
\texttt{std::priority\_queue<int, vector<int>, greater<int>>}, le tas min de la STL. Mêmes
valeurs pour les deux (fichiers triés de 200\,000 à 10\,000\,000 valeurs, puis mélangés).
Machine : Intel Core i5-1135G7, \texttt{g++} 15.2.

\begin{table}[H]
\centering
\small
\setlength{\tabcolsep}{4pt}
\begin{tabular}{|l|l|r|r|r|r|r|r|r|r|}
\hline
\multicolumn{2}{|c|}{} & \multicolumn{2}{c|}{\textbf{Insérer les $n$}} & \multicolumn{2}{c|}{\textbf{Extraire le min}}
  & \multicolumn{2}{c|}{\textbf{Insérer : reste}} & \multicolumn{2}{c|}{\textbf{Insérer : jusqu'à}} \\
\multicolumn{2}{|c|}{} & \multicolumn{2}{c|}{\textbf{valeurs} (ms)} & \multicolumn{2}{c|}{($\mu$s)}
  & \multicolumn{2}{c|}{\textbf{à la fin} ($\mu$s)} & \multicolumn{2}{c|}{\textbf{la racine} ($\mu$s)} \\
\hline
$n$ & \textbf{Ordre} & \texttt{TasMin} & \texttt{PQ} & \texttt{TasMin} & \texttt{PQ}
    & \texttt{TasMin} & \texttt{PQ} & \texttt{TasMin} & \texttt{PQ} \\
\hline
\multirow{2}{*}{200\,000}     & trié    & 2,0  & 1,95 & 0,98 & 1,15 & 0,09 & 0,08 & 0,34 & 0,30 \\
\cline{2-10}
                              & mélangé & \textbf{6,9} & 7,5 & \textbf{0,38} & 0,44 & 0,08 & 0,07 & \textbf{0,13} & 0,15 \\
\hline
\multirow{2}{*}{1\,000\,000}  & trié    & 9,2  & 9,1  & \textbf{1,76} & 2,33 & 0,07 & 0,07 & 0,59 & 0,53 \\
\cline{2-10}
                              & mélangé & \textbf{33} & 35 & 0,45 & 0,43 & 0,08 & 0,07 & \textbf{0,15} & 0,20 \\
\hline
\multirow{2}{*}{10\,000\,000} & trié    & 127  & \textbf{125} & 3,55 & \textbf{3,38} & 0,08 & 0,07 & 1,08 & 0,93 \\
\cline{2-10}
                              & mélangé & \textbf{361} & 394 & 3,51 & 3,52 & 0,08 & 0,08 & \textbf{0,20} & 0,21 \\
\hline
\multicolumn{2}{|l|}{\textit{Échanges de \texttt{TasMin}}} & \multicolumn{2}{c|}{trié : \textbf{0}} & \multicolumn{2}{c|}{17 à 23}
  & \multicolumn{2}{c|}{\textbf{0}} & \multicolumn{2}{c|}{17 à 23 ($\log_2 n$)} \\
\hline
\end{tabular}
\caption{Compilé avec l'optimiseur du compilateur (\texttt{-O3}). Médianes sur 15 exécutions ;
en gras : gagne au moins 11 fois sur 15.}
\end{table}

\begin{table}[H]
\centering
\small
\begin{tabular}{|l|r|r|r|r|}
\hline
& \multicolumn{2}{c|}{\textbf{Sans optimisation} (\texttt{-O0})} & \multicolumn{2}{c|}{\textbf{Optimisé} (\texttt{-O3})} \\
\hline
\textbf{Mesure} (1\,000\,000 valeurs) & \texttt{TasMin} & \texttt{PQ} & \texttt{TasMin} & \texttt{PQ} \\
\hline
Insérer les $n$ valeurs, trié (ms)           & \textbf{56,9} & 66,9          & 10,2          & \textbf{9,9} \\
\hline
Insérer les $n$ valeurs, mélangé (ms)        & 131           & \textbf{110}  & \textbf{34,3} & 37,8 \\
\hline
Extraire le min, mélangé ($\mu$s)            & 2,37          & \textbf{1,98} & 0,49          & 0,50 \\
\hline
Insérer : jusqu'à la racine, mélangé ($\mu$s)& 1,03          & \textbf{0,70} & \textbf{0,15} & 0,20 \\
\hline
Insérer : reste à la fin, mélangé ($\mu$s)   & \textbf{0,13} & 0,33          & 0,08          & \textbf{0,08} \\
\hline
\end{tabular}
\caption{Le même programme 13 compilé sans option (\texttt{g++ fichier.cpp}) puis avec
\texttt{-O3}. Médianes sur 7 exécutions ; en gras : gagne au moins 6 fois sur 7.}
\end{table}

\textbf{Ce qu'on en retient.}
\begin{enumerate}
  \item \textbf{Le fichier trié est le meilleur cas du tas min} : 0 échange, et l'insertion
        est 3 fois plus rapide que sur les valeurs mélangées (127 contre 361~ms à
        10\,000\,000). C'est l'inverse du tas max, pour lequel il était le pire cas.
  \item \textbf{Optimisé (\texttt{-O3}) : notre tas gagne l'insertion des valeurs mélangées}
        (7 à 9\,\% plus vite, 14 ou 15 fois sur 15). Sur le fichier trié, où il n'y a aucun
        échange, les deux versions sont presque égales (\texttt{PQ} 2\,\% devant à
        10\,000\,000).
  \item \textbf{Sans optimisation : \texttt{priority\_queue} gagne} dès qu'il faut faire
        remonter ou descendre une valeur (insertion mélangée 1,2 fois plus vite, extraction,
        insertion jusqu'à la racine) : elle déplace un \og trou \fg{} au lieu de faire un
        \texttt{swap}. Notre tas ne gagne que quand il n'y a rien à échanger (fichier trié,
        insertion qui reste à la fin). On ne compare donc deux versions qu'avec les
        \textbf{mêmes options de compilation}.
  \item \textbf{$O(\log n)$} : une insertion qui reste à la fin fait 0 échange, une qui monte
        jusqu'à la racine en fait 17 à 23, comme l'extraction, quand $n$ passe de 200\,000 à
        10\,000\,000.
\end{enumerate}

% ==============================================================================
\section{Conclusion}
% ==============================================================================

Le tas min est le tas max avec la comparaison inversée : même tableau, mêmes coûts
($\log_2 n$ échanges au plus pour insérer ou extraire le min). Mais l'ordre des données
s'inverse aussi : des valeurs croissantes, pire cas du tas max, sont le meilleur cas du tas
min (aucun échange). Face à \texttt{std::priority\_queue} avec \texttt{greater<int>}, le plus
rapide dépend encore de l'optimiseur du compilateur : \texttt{priority\_queue} sans
optimisation, notre tas avec \texttt{-O3} sur les valeurs mélangées.

\end{document}
